\newpage
% Оформляем заголовок аналогично СОДЕРЖАНИЮ, но без номера главы
\begin{center}
    {\bfseries ПЕРЕЧЕНЬ ГРАФИЧЕСКОГО МАТЕРИАЛА}
\end{center}

\addcontentsline{toc}{chapter}{ПЕРЕЧЕНЬ ГРАФИЧЕСКОГО МАТЕРИАЛА}

\vspace{1cm}

\noindent В состав графического материала дипломного проекта входят следующие документы, выполненные на листах формата А1:

\vspace{0.5cm}

\begin{enumerate}
    \item \textbf{БГУИР ДП 1-53 01 07 001 Э3}. Схема электрическая принципиальная узла мониторинга. Представлена на 1 листе формата А1. Отражает электрические связи между микроконтроллером ESP32 и периферийными датчиками SHT31.
    
    \item \textbf{БГУИР ДП 1-53 01 07 002 СД}. Схема данных (алгоритм функционирования). Представлена на 1 листе формата А1. Содержит логику работы встроенного ПО, включая циклы глубокого сна и алгоритмы фильтрации данных.
    
    \item \textbf{БГУИР ДП 1-53 01 07 003 ТЭП}. Технико-экономические показатели проекта. Представлена на 1 листе формата А1. Содержит диаграммы структуры затрат, график окупаемости и сравнительную таблицу характеристик.
    
    \item \textbf{БГУИР ДП 1-53 01 07 004 СМ}. Схема мониторинга и визуализации данных. Представлена на 1 листе формата А1. Демонстрирует архитектуру облачного сервиса и интерфейс пользователя (Dashboard).
\end{enumerate}

\vspace{1cm}

\noindent Все чертежи выполнены в соответствии с требованиями стандартов ЕСКД и ЕСПД. Оригиналы графического материала хранятся в отдельной папке и предъявляются Государственной экзаменационной комиссии в день защиты.

\vfill
\begin{center}
    Минск 2026
\end{center}

\newpage